\begin{textsample}{POS Dim 1 – sp – Score 29.00 – sp/sp\_ef\_ciencias\_humanas\_geo\_1.txt}  \label{ex:f1_pos_163}
Geografia A Base Nacional Comum Curricular ( BNCC ) estabelece para o componente de Geografia os conhecimentos, as competências e as habilidades \textbf{que} se espera \textbf{que} os estudantes desenvolvam no decorrer do Ensino Fundamental, e os propósitos \textbf{que} direcionam a educação brasileira para a formação humana integral e para a construção de \textbf{uma} sociedade \textbf{justa}, democrática e inclusiva.
O contato intencional e orientado com os conhecimentos geográficos é \textbf{uma} oportunidade para compreender o mundo em \textbf{que} se \textbf{vive}, na medida em \textbf{que} esse componente curricular aborda as ações humanas construídas nas distintas sociedades existentes nas diversas regiões do planeta.
Para fazer a leitura do mundo em \textbf{que} \textbf{vivem}, com base nas aprendizagens em Geografia, os estudantes precisam ser estimulados a pensar espacialmente, desenvolvendo o \textbf{raciocínio} geográfico.
Na Educação Básica, a Geografia permite ao estudante ler e interpretar o espaço geográfico por meio das formas, dos processos, das dinâmicas e dos fenômenos e a entender as relações entre as sociedades e a natureza em \textbf{um} mundo complexo e em constante transformação. > [... ] a Geografia, entendida como \textbf{uma} ciência social, \textbf{que} estuda o espaço construído pelo \textbf{homem}, a partir das relações \textbf{que} estes mantêm entre si e com a natureza, quer \textbf{dizer}, as questões da sociedade, com \textbf{uma} “ visão espacial ”, é por \textbf{excelência} \textbf{uma} \textbf{disciplina} formativa, capaz de instrumentalizar o aluno para \textbf{que} exerça de fato a sua cidadania. [... ] \textbf{Um} cidadão \textbf{que} reconheça o mundo em \textbf{que} \textbf{vive}, \textbf{que} se compreenda como indivíduo social capaz de construir a sua história, a sua sociedade, o seu espaço, e \textbf{que} \textbf{consiga} ter os mecanismos e os instrumentos para tanto. ( CALLAI, 2001, p.134 ) É importante reconhecer \textbf{que} o ensino de Geografia passou por \textbf{crises} e renovações.
As tensões, contradições e inspirações advindas de diferentes concepções do pensamento geográfico, por meio da Geografia Clássica ou Tradicional, a Geografia Neopositivista - ou Positivismo Lógico ou Geografia Teórico-Quantitativa -, a Geografia Crítica e a Geografia Humanista e Cultural, entre outras, contribuíram para a consolidação da Geografia Escolar, refletindo -se no processo de \textbf{ensino-aprendizagem} e na construção de políticas públicas educacionais.
Dessa forma, no ensino de Geografia, observa -se \textbf{uma} \textbf{expressiva} pluralidade de concepções \textbf{teórico-metodológicas} \textbf{que} orientam a prática docente e fundamentam a elaboração de propostas curriculares.
As transformações observadas apresentam pontos importantes para a reflexão sobre os conteúdos, as metodologias e as estratégias de avaliação e, sobretudo os caminhos para superar a dicotomia \textbf{historicamente} construída entre a Geografia Física e a Humana, \textbf{que} ainda persiste nos dias \textbf{atuais}, nas universidades e especialmente na Educação Básica.
No entanto, apesar do reconhecimento das diferentes contribuições, o Currículo Paulista apresenta temáticas e \textbf{abordagens} próximas da Geografia Crítica, Humanista e Cultural, quando se opta por enfatizar a relação sociedade e natureza e a necessidade de se refletir, agir e fazer escolhas sustentáveis diante dos desafios \textbf{contemporâneos}.
O Currículo Paulista de Geografia do Ensino Fundamental está organizado com base nos princípios e conceitos da Geografia \textbf{contemporânea}.
Ressalta -se \textbf{que}, embora o espaço seja o conceito mais amplo e complexo da Geografia, é necessário \textbf{que} os estudantes dominem outros conceitos operacionais, \textbf{que} expressam aspectos diferentes do espaço geográfico : território, lugar, região, natureza e paisagem.
Diante da complexidade do espaço geográfico, o ensino de Geografia, na \textbf{contemporaneidade}, tem o desafio de articular \textbf{teorias}, pressupostos éticos e políticos da educação, bem como caminhos \textbf{metodológicos} ; para \textbf{que} os estudantes aprendam a pensar e a reconhecer o espaço por meio de diferentes escalas e tempos, desenvolvendo \textbf{raciocínios} geográficos, o pensamento espacial e construindo novos conhecimentos. > Pensar espacialmente, compreendendo os conteúdos e conceitos geográficos e suas representações, também envolve o \textbf{raciocínio}, definido pelas habilidades \textbf{que} desenvolvemos para compreender, a estrutura e a função de \textbf{um} espaço e descrever sua organização e relação a outros espaços, portanto, analisar a ordem, a relação e o padrão dos objetos espaciais. ( CASTELLAR, 2017, p.164 ) O \textbf{raciocínio} geográfico está relacionado com \textbf{uma} maneira de exercitar o pensamento espacial, por meio de princípios fundamentais : - Analogia : \textbf{um} fenômeno geográfico sempre é comparável a outros.
A identificação das semelhanças entre fenômenos geográficos é o início da compreensão da unidade terrestre ; - Conexão : \textbf{um} fenômeno geográfico nunca acontece isoladamente, mas sempre em interação com outros fenômenos próximos ou distantes ; - Diferenciação : é a variação dos fenômenos de interesse da geografia pela superfície terrestre ( por exemplo, o clima ), resultando na diferença entre áreas ; - Distribuição : exprime como os objetos se repartem pelo espaço ; - Extensão : espaço finito e contínuo delimitado pela ocorrência do fenômeno geográfico ; - Localização : posição particular de \textbf{um} objeto na superfície terrestre.
A localização pode ser absoluta ( definida por \textbf{um} sistema de coordenadas geográficas ) ou relativa ( expressa por meio de relações espaciais topológicas ou por interações espaciais ) ; - Ordem : ordem ou arranjo espacial é o princípio geográfico de maior complexidade.
Refere -se ao modo de estruturação do espaço de acordo com as regras da própria sociedade \textbf{que} o produziu.
O ensino de Geografia mobiliza competências e habilidades por meio de diferentes linguagens, de princípios e dos conceitos estruturantes espaço geográfico, paisagem, lugar, território e região e outras \textbf{categorias} \textbf{que} contemplam a natureza, a sociedade, o tempo, a cultura, o trabalho e as redes, entre outros, considerando as suas diversas escalas.
Outro conceito estruturante refere -se à educação \textbf{cartográfica}, \textbf{que} deve perpassar todos os anos do Ensino Fundamental.
Quanto às \textbf{categorias}, especialmente no \textbf{que} se refere à natureza e sociedade, é necessário aprofundar o estudo sobre os \textbf{fundamentos} do pensamento científico e filosófico. > Para entender o ensino, a prática do ensino de Geografia, é \textbf{preciso} pensar, pois, nas bases da ciência de referência.
Na atualidade, a ciência geográfica tem passado por algumas mudanças.
A Geografia é \textbf{um} campo do conhecimento científico multidimensional, sempre buscou compreender as relações \textbf{que} se estabelecem entre o \textbf{homem} e a natureza e como essas relações vêm constituindo diferentes espaços ao longo da história. \textbf{Hoje}, mais do \textbf{que} nunca, essa busca leva ao surgimento de \textbf{uma} pluralidade de caminhos.
As relações sociais, as práticas sociais geram e são geradas por espacialidades complexas, \textbf{que} \textbf{demandam} diferentes olhares, ampliando consideravelmente o campo temático e os problemas tratados pela Geografia.
E o ensino dessa \textbf{disciplina}, o \textbf{que} tem a \textbf{ver} com essa realidade?
As preocupações \textbf{que} orientam a produção científica da Geografia no âmbito acadêmico são as mesmas \textbf{que} norteiam a estruturação da \textbf{disciplina} escolar?
Sim e não.
Sim, porque as duas têm a mesma base epistemológica ; não, porque na escola existem \textbf{influências} diversas \textbf{que} dão \textbf{um} contorno peculiar a essa área do conhecimento.
O \textbf{que} valida a geografia escolar é a sua base, sua ciência de referência. ( CAVALCANTI, 2012, p.90 )

% matched lemmas: abordagem, atual, cartográfico, categoria, conseguir, contemporaneidade, contemporâneo, crise, demandar, disciplina, dizer, ensino-aprendizagem, excelência, expressivo, fundamento, historicamente, hoje, homem, influência, justo, metodológico, preciso, que, raciocínio, teoria, teórico-metodológico, um, ver, viver
\end{textsample}
